% !TeX root = ../thesis.tex
\abstract
\begin{singlespace}
\noindent This study examined how to transfer clinical records to a research centre while preserving the measurements needed for analysis. A workstation sent data to a server at the Kwame Nkrumah University of Science and Technology, Kumasi. The public PhysioNet sepsis cohort contained 790,215 hourly records from 20,336 patients. Feature selection, principal component analysis (PCA), autoencoders and lossless formats were compared on size, preparation and delivery time, and preservation of observed values. At an 80\% hourly missingness threshold, filtering alone retained 13 of 40 predictors, while protection rules retained 28. Parquet reduced the source CSV from 75.69 to 9.17 MB. At 16 components, the autoencoder recovered more observed variation than PCA (\(R^2=0.907\) and \(0.676\), respectively), but required more time to fit. All 586 SCP package transfers passed remote checksum verification, and the receiver recovered all 82 package configurations. With equal weights across seven decision criteria, Parquet ranked first in all eight fresh-data groups that permitted approximation. In a separate paired experiment on a 42.91 MB archive, rsync completed verified delivery a median of 10.46 seconds sooner than SCP after a controlled interruption. On the route studied, Parquet offers a practical starting point for fresh tables. SCP was quicker for the two smaller prepared archives on uninterrupted transfers. For the large archive exposed to interruption, rsync with partial-file recovery is preferable.
\end{singlespace}
